%HEADER

\documentclass[11pt,a4paper]{scrreprt}	%artikeltyp

\usepackage[english]{babel}	%sprache
\usepackage[utf8]{inputenc}%utf8, für zeichen, umlaute etc.
\usepackage[a4paper,lmargin={2.5cm},rmargin={2.5cm},tmargin={2.5cm},bmargin = {2.5cm}]{geometry}	%Abstände, usw. ändern
\usepackage{amsmath}	%mathe
\usepackage{subfiles}
\usepackage{amssymb}	%mathe
%\usepackage{amsfonts}	%mathe
\usepackage{nicefrac}	%erlaubt z.B. schräge brüche(\nicefrac)
\usepackage{graphicx}% Graphikpaket (für PNG,GIF,JPG)
\usepackage{float,rotating}%Option für fixe Position; Drehung von Abbildungen etc.
\usepackage{setspace}%zeilenabstand 
\usepackage{hyperref}%klickbare links
\usepackage{gensymb}%z.B. Gradzeichen
\usepackage{wrapfig} %gibt wrapfigures
\usepackage{epstopdf} % eps dateien einbinden
\usepackage{dsfont}
\usepackage{subcaption}
\usepackage{MnSymbol}	%für Kreisbogenzeichen
\usepackage{subfiles}
\usepackage{tabularx}
\usepackage{mhchem}
\usepackage[round]{natbib}

%..............Kopf und Fußzeilengestaltung:.................
\usepackage{fancyhdr}
\pagestyle{headings}

%\lhead{\nouppercase{\leftmark}}
%\rhead{\thepage}



%..............Definitionen:................
\newcommand{\pd}{\partial}
\newcommand{\pfrac}[2]{\frac{\partial #1}{\partial #2}}
\newcommand{\ph}{\varphi}
\newcommand{\fem}{\frac{e}{m}}
\newcommand{\om}{\omega}
\newcommand{\und}{~\text{und}~}
\newcommand{\dg}{\degree}
\newcommand{\un}[1]{\,\mathrm{#1}}
\newcommand{\cd}{\cdot}
\renewcommand*{\chapterheadstartvskip}{\vspace*{0.5cm}}%abstand oben bei anfang eines chapters
\newcommand{\R}{\grqq}
\newcommand{\A}{\glqq}
%\renewcaptionname{ngerman}{\figurename}{Abb.} %"`Abb."' statt "`Abbildung"' gilt nur wenn ngerman geladen ist
\setlength\parindent{0pt}


\begin{document}

	\begin{titlepage}
		
		\vspace{10 cm}
		\begin{figure}
			\centering
			\includegraphics*[width=0.55\linewidth]{Bilder/Uni_logo.pdf}
		\end{figure}
	
		
		\begin{center}
			\rule{\textwidth}{1.6pt}\vspace*{-\baselineskip}\vspace*{2pt} % Thick horizontal rule
			\rule{\textwidth}{0.4pt} % Thin horizontal rule
			
			\vspace{1.5\baselineskip} % Whitespace above the title
			
			\textbf{\LARGE Reconstructing North Atlantic Ocean Heat Content Using Convolutional Neural Networks} % Title
			
			\vspace{0.75\baselineskip} % Whitespace below the title
			
			\rule{\textwidth}{0.4pt}\vspace*{-\baselineskip}\vspace{3.2pt} % Thin horizontal rule
			\rule{\textwidth}{1.6pt} % Thick horizontal rule
			
			\vspace{2\baselineskip} % Whitespace after the title block
			\vspace*{2.0cm}
			\LARGE{Master Thesis\\}
			\vspace{3cm}
			\Large 
			Author: Paul Simon Lentz\\
			Matriculation Number: 737 4279 \\
			University Hamburg \\
			Master of Integrated Climate System Sciences \\
			Faculty for Mathematics, Informatics and Natural Sciences \\
			Supervisors: Prof. Dr. Johanna Baehr, Dr. Christopher Kadow \\
			Deadline: 16.01.2023\\
			\vspace{5 cm}
		\end{center}
	
	
	\end{titlepage}
	
	\pagenumbering{gobble}
	
	%\addtocounter{page}{-5}
	\thispagestyle{empty}
	\clearpage\mbox{}\clearpage
	
	\onehalfspacing
	\tableofcontents %inhaltsverzeichnis
	\listoffigures
	\listoftables
	
	\newpage
	\pagenumbering{arabic}
	\setlength\parindent{0cm}
	
	\subfile{./Subfiles/Abstract.tex}
	
	\subfile{./Subfiles/Introduction.tex}
	
	\subfile{./Subfiles/Data_Methods.tex}
	
	\subfile{./Subfiles/Results.tex}
	
	%\subfile{Diskussion.tex}	
	
	\bibliographystyle{apalike}
	\bibliography{Masterbib}
	

	%\subfile{Anhang.tex}


\end{document}